\subsection{Circuito resistivo}

Inicialmente se analizó la configuración correspondiente al esquematico visto en la Figura~\ref{fig:tp3_resistencia}. Para ello, se mantuvo cerrada únicamente la llave $S_R$, de modo que la carga conectada estuvo formada por el banco de lámparas incandescentes, modelado mediante una resistencia equivalente $R$.

\begin{figure}[H]
    \centering
    \begin{circuitikz}
        \draw
        (0,0) to[sV,l=$V_S$] (0,4)
              -- (2.5,4)
              to[closing switch,l=$S_R$] (2.5,2.5)
              to[american resistor,l=$R$] (2.5,0)
              -- (0,0);
    \end{circuitikz}
    \caption{Configuración correspondiente al circuito resistivo.}
    \label{fig:tp3_resistencia}
\end{figure}

\subsubsection{Corroboración experimental de los resultados}

A partir de las mediciones realizadas se obtuvieron los valores presentados en el Tabla~\ref{tab:tp3_paso1}. En esta configuración se midieron la tensión eficaz aplicada sobre la carga, la corriente eficaz total consumida por el circuito y la potencia activa entregada a la rama resistiva. A partir de estas mediciones, se calcularon la potencia aparente, la potencia reactiva, el factor de potencia y el módulo de la impedancia equivalente.

\begin{table}[H]
    \centering
    \caption{Mediciones y magnitudes calculadas para el ensayo resistivo.}
    \label{tab:tp3_paso1}
    \begin{tabular}{lc}
        \toprule
        Variable & Valor \\
        \midrule
        $R$ [$\Omega$] & 166,7 \\
        $|\underline{V}|$ [V] & 100 \\
        $|\underline{I}|$ [A] & 0,60 \\
        $P$ [W] & 62 \\
        $S$ [VA] & 62 \\
        $Q$ [VAr] & $\approx 0$ \\
        $fp$ & $\approx 1,00$ \\
        \bottomrule
    \end{tabular}
\end{table}

\subsubsection{Triángulo de potencia}

A partir de los valores de Tabla~\ref{tab:tp3_paso1} se construyó el triángulo de potencia correspondiente. Como se observa en la Figura~\ref{fig:triangulo_resistivo}, la potencia reactiva resultó aproximadamente nula, por lo que la potencia aparente coincidió con la potencia activa. Este comportamiento era el esperado para una carga puramente resistiva.

\begin{figure}[H]
    \centering
    \triangulopot{62}{0}{62}{Solo $R$}
    \caption{Triángulo de potencia correspondiente al circuito resistivo.}
    \label{fig:triangulo_resistivo}
\end{figure}

\subsubsection{Analisis de los resultados}

Debido a que se trataba de un circuito compuesto únicamente por una resistencia y una fuente de tensión alterna, se espera que la potencia reactiva \(Q\) sea nula.

Además, al ser un circuito puramente resistivo, la tensión y la corriente debían encontrarse en fase, por lo que el factor de potencia esperado era \(1\). Sin embargo, al calcularlo experimentalmente se obtuvo un valor de \(1{,}033\), lo que representa un error aproximado del \(3{,}2\%\). Este resultado no es físicamente posible, ya que el factor de potencia no puede ser mayor que uno. Por lo tanto, como el factor de potencia se calculó mediante la relación en la eq
:\eqref{eq:factor_potencia} la diferencia se atribuyó a pequeñas variaciones en las lecturas y a la precisión limitada del voltímetro, amperímetro y vatímetro.

De hecho, se pudo identificar una inconsistencia en las mediciones, ya que, al tratarse de un circuito puramente resistivo, la potencia aparente calculada como en la eq:\eqref{eq:potencia_aparente}, debería coincidir con \(P\) medido con el vatímetro. Sin embargo, el valor de \(P\) resultó aproximadamente un \(3{,}2\%\) mayor que \(S\), lo cual explica que el factor de potencia calculado usando eq
:\eqref{eq:factor_potencia} haya dado un valor superior a la unidad.

Finalmente, al simular el circuito, la corriente obtenida fue de aproximadamente \(0{,}5997\,\mathrm{A}\), mientras que el valor real medido con el amperímetro fue de \(0{,}6\,\mathrm{A}\). Este resultado muestra una muy buena concordancia entre la simulación y la medición experimental. En este caso, el comportamiento no óhmico de las lámparas incandescentes no produjo una diferencia significativa, ya que los valores utilizados para calcular la resistencia de las lámparas fueron obtenidos en esta misma etapa de medición. Por lo tanto, las lámparas no llegaron a modificar considerablemente su resistencia por efecto del calentamiento del filamento.
